\documentclass[english, parskip=full]{scrreport}
\usepackage[utf8]{inputenc}  % Kodierung der Datei
\usepackage[english]{babel}  % Sprache auf Deutsch 
\usepackage{color}
\usepackage{graphicx, gensymb}
\usepackage{amsfonts, cancel, amssymb}
\usepackage{amsmath, amsthm, array, esint}
\usepackage{booktabs, multirow}  % schönere Tabellen
\usepackage{caption, slashed}  % Anpassung der Captions
\usepackage{scrlayer-scrpage}    % Manipulation des Seitenstils
\pagestyle{scrheadings}  % Stil für die Seite setzen
\automark{section}  % Abschnittsnamen als Seitenbeschriftung 
\setheadsepline{.5pt}    % Kopzeile durch Linie abtrennen
\usepackage[hidelinks]{hyperref}  % Links und weitere PDF-Features
\usepackage{typearea, tikz, pgffor, verbatim}
\usetikzlibrary{calc, quotes}
\usepackage[cache=false, outputdir=build]{minted}
\usepackage{placeins} % provides \FloatBarrier

\definecolor{bg}{rgb}{0.9,0.9,0.9}
\newcommand\numberthis{\addtocounter{equation}{1}\tag{\theequation}}
\newcommand{\bra}{}
\newcommand{\kom}{}
\newcommand{\brak}{}
\newcommand{\braket}{}
\newcommand{\intx}{}
\newcommand{\intr}{}
\newcommand{\kla}{}
\newcommand{\klam}{}
\newcommand{\klammer}{}
\newcommand{\oper}{}
\newcommand{\tecxt}{}
\newcommand{\Bra}{}
\newcommand{\Ket}{}
\renewcommand{\i}{\text{i}}
\newcommand{\e}{\text{e}}
\newcommand*{\Fourier}{\textsc{Fourier}\ }
\newcommand*{\F}{\mathcal{F}}
\newcommand*{\sinc}{\text{sinc\,}}
\newcommand{\conv}{\mathop{\scalebox{3.5}{\raisebox{-0.38ex}{\makebox[0.5\width][c]{$\ast$}}}}\limits}

\newcommand*{\titel}{Zusammenfassung der Vorlesungen}
\newcommand*{\autor}{Joachim Schwardt}

\hypersetup{pdfauthor={\autor}, pdftitle={\titel}}

%\titlehead{titlehead}
\subject{Theoretische Physik}
\title{\titel}
\author{\autor}
\date{\today}

\begin{document}

%\begin{titlepage}
\maketitle  % Titel setzen
%\tableofcontents  % Inhaltsverzeichnis setzen
%\end{titlepage}


\chapter{Quantentheorie II}
\section{Relativistische Quantenmechanik}
\subsection{Ortsraum und Teilchen in einer Dimension}
Der Translationsoperator ist definiert als
\begin{align}
\hat{T}_a &:= \e^{-\i\hat{p}a}.
\end{align}
Es gilt für den Kommutator mit dem Ortsoperator dann
\begin{align}
\klam\left[ \hat{x}, \hat{T}_a \right] &= \sum_{n\ge 0}\frac{(-\i a)^n}{n!}\klam\left[ \hat{x}, \hat{p}^n \right] = \sum_{n\ge 0}\frac{(-\i a)^n}{n!} n\i\hat{p}^{n-1} = a\hat{T}_a.
\end{align}
Der Ortseigenwert von einem Zustand $\hat{T}_a\Ket | x_0 \rangle$ folgt dann zu
\begin{align}
\hat{x}\kla\left( \hat{T}_a\Ket | x_0 \rangle \right) &= \hat{T}_a\kla\left( x + a \right)\Ket | x_0 \rangle.
\end{align}
Es gilt also $\hat{T}_a\Ket | x_0 \rangle\propto \Ket | x_0 + a \rangle$.
Da $\hat{T}_a^\dagger\hat{T}_a = 1$ gilt insbesondere die Gleichheit, also
\begin{align}
\hat{T}_a\Ket | x \rangle &= \Ket | x + a \rangle.
\end{align}
Insbesondere verwenden wir die Konvention $\brak\langle p | x \rangle = \e^{-\i px}$.

%Insbesondere können wir $\hat{T}_a$ im folgenden Skalarprodukt auch ``nach links'' anwenden, also
%\begin{align}
%\brak\langle p | x + a \rangle &= \braket\langle p | \hat{T}_a | x \rangle = \e^{-\i pa}\brak\langle p | x \rangle.
%\end{align}
%Andererseits ist dann
%\begin{align}
%\brak\langle p | x \rangle &= \intr\int_{\mathbb{R}^{ }} \brak\langle p | y \rangle\brak\langle y | x \rangle \, \mathrm{d}^{ }y = \intr\int_{\mathbb{R}^{ }} \e^{-\i py}\brak\langle p | 0 \rangle \delta(x-y) \, \mathrm{d}^{ }y
%\end{align}


\subsection{Rotationsinvarianz und Drehimpuls}
Rotationen werden im $\mathbb{R}^3$ durch die Rotationsmatrizen
\begin{align}
R_{\vec{n}} (\theta) 
\end{align}
mit Einheitsvektor $\vec{n}$ beschrieben.
Für $\vec{n} = \vec{e}_z$ ist die Drehung dann beispielsweise gegeben durch
\begin{align}
R_{\vec{n}} (\theta) &:= R_z(\theta) = \begin{pmatrix}
\cos\theta & -\sin\theta & 0\\
\sin\theta & \cos\theta & 0\\
0 & 0 & 1
\end{pmatrix}.
\end{align}
Die zugehörige infinitesimale Drehung ist
\begin{align}
R_z(\varepsilon) &= 1 -\i\varepsilon\begin{pmatrix}
0 & -\i & 0 \\ \i & 0 & 0 \\ 0&0&0
\end{pmatrix} + \mathcal{O}(\varepsilon^2) \equiv 1 - \i\varepsilon\ell_z + \mathcal{O}(\varepsilon^2).
\end{align}
Die Matrizen $(\ell_m)_{jk} = -\i\epsilon_{jkm}$ heißen Generatoren und spannen eine Lie-Algebra auf.
Ihre explizite Form ist
\begin{align}
\ell_x &= \begin{pmatrix}
0 & 0 & 0 \\ 0 & 0 & -\i \\ 0&\i&0
\end{pmatrix}, &&& \ell_y &= \begin{pmatrix}
0 & 0 & \i \\ 0 & 0 & 0 \\ -\i&0&0
\end{pmatrix} &&\tecxt\text{und} & \ell_z &= \begin{pmatrix}
0 & -\i & 0 \\ \i & 0 & 0 \\ 0&0&0
\end{pmatrix}.
\end{align}
Mit der Relation $\epsilon_{abc}\epsilon_{ade} = \delta_{bd}\delta_{ce} - \delta_{be}\delta_{cd}$ kann man zeigen, dass
\begin{align}
 \ell_n\ell_m  &\equiv  -\epsilon_{akn}\epsilon_{kcm} = \epsilon_{kan}\epsilon_{kcm} = \delta_{ac}\delta_{nm} - \delta_{am}\delta_{nc}
\end{align}
und somit
\begin{align}
\klam\left[ \ell_n,\ell_m \right] &= \ell_n\ell_m - \ell_m\ell_n \equiv \delta_{ac}\delta_{nm} - \delta_{am}\delta_{nc} - \delta_{ac}\delta_{mn} + \delta_{an}\delta_{mc} = \\
&= \delta_{an}\delta_{mc} - \delta_{am}\delta_{nc} = \underbrace{\epsilon_{kac}}_{=\,(\ell_k)_{ac}} \epsilon_{knm} \equiv \i\epsilon_{knm}\ell_k.
\end{align}
Es gilt außerdem
\begin{align}
R_{\vec{n}}(\theta) &= \e^{-\i\theta\vec{n}\cdot\vec{\ell}} \equiv \e^{-n_m(\ell_m)_{jk}}.
\end{align}
Eine Abbildung $R\mapsto D(R)$ mit $D(R_1R_2)=D(R_1)D(R_2)$ heißt Darstellung einer Drehung.
Beispielsweise gelten
\begin{align}
\Ket | \psi \rangle &\rightarrow \Ket | \psi' \rangle = D(R)\Ket | \psi \rangle, &&& Q&\rightarrow Q'=Q, &&&\vec{p}&\rightarrow \vec{p}\,'= D(R)\vec{p}.
\end{align}
Im Folgenden werden die Korrekturen $\mathcal{O}(\varepsilon^2)$ nicht mehr explizit geschrieben.
Die Transformation für eine Ladungsverteilung führt dann auf
\begin{align}
\rho(\vec{r}) &\rightarrow \rho'(\vec{r}\,') = \rho(R_z\vec{r}) = \rho(x-\varepsilon y, y + \epsilon x, z) = \\
&= \rho(\vec{r}) - \epsilon\partial_y\rho(\vec{r}) + \epsilon\partial_x\rho(\vec{r}) = \kla\left( 1 - \i\varepsilon\klam\left[ x\partial_y - y\partial_x \right] \right)\rho(\vec{r}).
\end{align}
Wir definieren
\begin{align}
J_m &:= \i\epsilon_{mjk}x_j\partial_{x_k}.
\end{align}
Für eine beliebige Quantentheorie fordern wir 
\begin{enumerate}
\item Symmetrie, also $\hat{D}$ unitär damit
\begin{align}
\brak\langle \psi' | \phi' \rangle &= \braket\langle \psi | \hat{D}^\dagger\hat{D} | \phi \rangle \overset{!}{=} \brak\langle \psi | \phi \rangle.
\end{align}
und zudem $\klam\left[ \hat{H},\hat{D} \right] \overset{!}{=} 0$.
\item Infinitesimale Repräsentation über
\begin{align}
\hat{D}_m &= 1 - \i\varepsilon\hat{J}_m
\end{align}
mit $\hat{J}_m^\dagger = \hat{J}_m$ hermitisch, $\klam\left[ \hat{J}_j, \hat{J}_k \right] = \i\epsilon_{jkm}\hat{J}_m$ und $\klam\left[ \hat{H}, \hat{J}_m \right] \overset{!}{=} 0$.
\end{enumerate}
Insbesondere gilt dann
\begin{align}
\klam\left[ \hat{J}^2, \hat{J}_m \right] &= \klam\left[ \hat{J}_j\hat{J}_j, \hat{J}_m \right] = \hat{J}_j\i\epsilon_{jkm}\hat{J}_k + \i\epsilon_{jkm}\hat{J}_k \hat{J}_j = \\
&= \i\epsilon_{jkm}\hat{J}_j\hat{J}_m + \i\epsilon_{kjm}\hat{J}_j\hat{J}_m = \i\hat{J}_j\hat{J}_m\kla\left( \epsilon_{jkm} + \epsilon_{kjm} \right) = 0.
\end{align}
Ohne Beschränkung können wir eine Raumrichtung auszeichnen, indem wir $\hat{J}^2$ und $\hat{J}_z$ gemeinsam diagonalisieren.
Sei also $\hat{J}^2\Ket | j,m \rangle = j(j+1)\Ket | j,m \rangle$ und $\hat{J}_z\Ket | j,m \rangle = m\Ket | j,m \rangle$.
Die ausführliche Konstruktion von Leiteroperatoren $\hat{J}_\pm := \hat{J}_x\pm\i\hat{J}_y$ wurde in Quantentheorie I besprochen.
Man findet letztlich $j\in\klammer\left\{ 0,\frac{1}{2},1,\frac{3}{2},\dots \right\}$ und $|m|\le j$.
\\
Die einfachste nicht-triviale Darstellung hat daher $j=\frac{1}{2}$, was auf die Pauli-Matrizen führt.
Definiere dazu die zwei Zustände $\Ket | \pm \rangle := \Ket | j=\frac{1}{2}, m=\pm\frac{1}{2} \rangle$, womit wir einen beliebigen Zustand gemäß
\begin{align}
\Ket | \psi \rangle &= \psi_+\Ket | + \rangle + \psi_- \Ket | - \rangle
\end{align}
zerlegen können.
Für den Vektor $\psi := (\psi_+,\psi_-)^\top$ gilt dann
\begin{align}
\psi&\rightarrow \psi' = \kla\left( 1 - \i\theta\frac{\sigma_k}{2} \right)\psi,
\end{align}
wobei die Pauli-Matrizen $\sigma_k$ die Generatoren der Gruppe SU(2) sind, welche isomorph zu SO(3) ist (also insbesondere gleiche Kommutatorrelationen).

\subsection{Lorentzinvarianz}
Wir verwenden hier die Vorzeichenkonvention $g_{\mu\nu} = (+---)$.
Die allgemeine Lorentztransformation wird durch $\Lambda^\mu_{\ \nu}$ beschrieben, wobei
\begin{align}
x^\mu&\rightarrow x'^\mu = \Lambda^\mu_{\ \nu} x^\nu.
\end{align}
Das Skalarprodukt bleibt hierbei erhalten, das heißt $x'^\mu y'_\mu = x^\mu y_\mu$.
Die infinitesimale Beschreibung ist
\begin{align}
\Lambda^\mu_{\ \nu} (\varepsilon) &= 1 - \i\epsilon G,
\end{align}
wobei es hier sechs Generatoren $G$ gibt.
Diese teilen sich auf in drei Boosts
\begin{align}
k_x &= \begin{pmatrix}
0&\i&0&0 \\ \i&0&0&0 \\ 0&0&0&0 \\ 0&0&0&0
\end{pmatrix}, &&& k_y &= \begin{pmatrix}
0&0&\i&0 \\ 0&0&0&0 \\ \i&0&0&0 \\ 0&0&0&0
\end{pmatrix}, &&& k_z &= \begin{pmatrix}
0&0&0&\i \\ 0&0&0&0 \\ 0&0&0&0 \\ \i&0&0&0
\end{pmatrix}, 
\end{align}
und drei Rotationen
\begin{align}
\ell_x &= \begin{pmatrix}
0&0&0&0 \\ 0&0&0&0 \\ 0&0&0&-\i \\ 0&0&\i&0
\end{pmatrix}, &&& \ell_y &= \begin{pmatrix}
0&0&0&0 \\ 0&0&\i&0 \\ 0&0&0&0 \\ -\i&0&0&0
\end{pmatrix}, &&& \ell_z &= \begin{pmatrix}
0&0&0&0 \\ 0&0&-\i&0 \\ 0&\i&0&0 \\ 0&0&0&0
\end{pmatrix}.
\end{align}
Es gelten die Kommutatorrelationen
\begin{align}
\klam\left[ \ell_x, \ell_y \right] &= \i\ell_z , &&& \klam\left[ k_x,k_y \right] &= -\i\ell_z &&\tecxt\text{und} & \klam\left[ \ell_x, k_y \right] &= \i k_z &&\tecxt\text{(zyklisch)} .
\end{align}
Man führt zudem $K_m, J_m$ für die Darstellung ein, welche dieselben Relationen 
\begin{align}
\klam\left[ J_x, J_y \right] &= \i J_z , &&& \klam\left[ K_x, K_y \right] &= -\i J_z &&\tecxt\text{und} & \klam\left[ J_x, K_y \right] &= \i K_z &&\tecxt\text{(zyklisch)} .
\end{align}
erfüllen.

\subsubsection{Darstellung für Spin $\frac{1}{2}$}
Eine mögliche Darstellung ist $J_m := \frac{\sigma_m}{2}$ und $K_m := \mp\i\frac{\sigma_m}{2}$.
Man nennt dies auch die zweidimensionale komplexe Darstellung der Lorentz-Algebra.
Unter Raumspiegelung transformieren sich die folgenden Größen wie
\begin{align}
x&\rightarrow -x, &&& p&\rightarrow -p, &&& J&\rightarrow J && \tecxt\text{und} & K&\rightarrow -K.
\end{align}
TODO: Dirac-Spinoren, $\gamma$-Matrizen, Weyl-Spinoren, Ausblick auf weitere Darstellungen der L-Algebra.

\subsection{Relativistische Wellengleichung}
Aus der Elektrodynamik kennen wir beispielsweise $\square\phi = 0$ und $(\nabla^2 + \vec{k}^2)\phi = 0$ oder $\i\partial_t + \frac{1}{2m}\nabla^2\phi = 0$ aus Quantentheorie I.
Diese Gleichungen sind alle rotationsinvariant, weshalb der Operator $\nabla^2$ auftaucht.
Frage: Welche Gleichungen erlaubt die Lorentz-Invarianz?
\\
\subsubsection{Klein-Gordon-Gleichung}
Die einfachste relativistische Gleichung ist die Klein-Gordon-Gleichung als Verallgemeinerung der Schrödingergleichung für ein Skalarfeld mit $\phi'(x') = \phi(x)$.
Sie lautet 
\begin{align}
\kla\left( \partial_\mu\partial^\mu + m^2 \right)\phi &= 0.
\end{align}
Wir betrachten den nicht-relativistischen Limes, also $\psi =: \e^{-\i mt}\psi_\tecxt\text{nr} $ mit $E=m+\mathcal{O}(\varepsilon)$.
Dann ist
\begin{align}
0 &= \kla\left( -m^2 - \nabla^2 - 2\i m\partial_t + \mathcal{O}\kla\left( \partial_t^2\psi_\tecxt\text{nr} \right) + m^2 \right)\psi_\tecxt\text{nr},
\end{align}
was auf die Schrödingergleichung
\begin{align}
\i\partial_t\psi_\tecxt\text{nr} &= -\frac{1}{2m}\nabla^2\psi_\tecxt\text{nr}
\end{align}
führt, wenn wir die höheren Ordnungen der Zeitableitung vernachlässigen.
Die KG-Gleichung findet Anwendung in der QFT für Feldoperatoren $\hat{\phi}(x)$.
Die Lagrange-Dichte lautet
\begin{align}
\mathcal{L}_\tecxt\text{KG} &= \kla\left( \partial_\mu\phi^\ast \right)\kla\left( \partial^\mu \phi \right) - m^2\phi^\ast\phi
\end{align}
und führt mit 
\begin{align}
\frac{\partial\mathcal{L}}{\partial \phi} &= \partial^\mu \frac{\partial\mathcal{L}}{\partial (\partial^\mu\phi)}
\end{align}
wieder auf die Klein-Gordon-Gleichung (analog mit $\phi^\ast$).
Für die Schrödingergleichung konnten wir $\rho = \phi^\ast\phi$ als eine Wahrscheinlichkeitsdichte und 
\begin{align}
\vec{j} &= -\frac{\i}{2m}\kla\left( \phi^\ast\nabla\phi - \phi\nabla\phi^\ast \right)
\end{align}
als eine Stromdichte für $\rho$ interpretieren.
Es galt eine Kontinuitätsgleichung der üblichen Form $\partial_t\rho + \nabla\cdot\vec{j} = 0$.
Für Klein-Gordon findet man ähnlich
\begin{align}
0 &= \phi^\ast\kla\left( \square + m^2 \right)\phi -\phi\kla\left( \square + m^2 \right)\phi^\ast = \partial_\mu \kla\left( \phi^\ast\partial^\mu\phi - \phi\partial^\mu\phi^\ast \right),
\end{align}
also $\partial_\mu j^\mu = 0$ für die Viererstromdichte %$j^\mu = (\rho, \vec{j})$ mit 
\begin{align}
j^\mu := \frac{1}{2m}\kla\left( \phi^\ast \partial^\mu\phi - \phi \partial^\mu\phi^\ast \right).
\end{align}
Allerdings ist hier $\rho = j^0 < 0$ möglich, wir können dies also nicht mehr als Wahrscheinlichkeitsdichte interpretieren.
Für die zwei ebenen Wellen $\phi_\pm \propto \e^{\mp\i(Et - \vec{p}\cdot\vec{r})}$ ist vielmehr $\rho = \pm \frac{E}{m}$.
Wir finden also Anti-Teilchen mit $E<0$ und $\rho< 0$.
Diese Interpretation kann zwar funktionieren, aber wir haben jetzt zwei Teilchenarten.
Eigentlich sollte KG aber eine Einteilchentheorie beschreiben (Ausblick: QFT).
\subsubsection{Dirac-Gleichung}
Idee: $E=\sqrt{m^2 + \vec{p}\,^2}\overset{!}{=} \beta m + \vec{\alpha}\cdot\vec{p}$, also
%\begin{align}
%\i\partial_t &\equiv \sqrt{m^2 - \nabla^2}\overset{!}{=}: \beta m - \vec{\alpha}\cdot\nabla.
%\end{align}
\begin{align}
E^2 &= \beta^2m^2 + \kla\left( \vec{\alpha}\cdot\vec{p} \right)^2 + \klammer\left\{ \beta, \vec{\alpha} \right\}\cdot m\vec{p} + \klammer\left\{ \alpha_j, \alpha_k \right\}p_jp_k \overset{!}{=} m^2 + \vec{p}\,^2.
\end{align}
Also müssen $\beta^2 = 1$, $\alpha_j^2 = 1$, $\klammer\left\{ \beta, \alpha_j \right\} = 0$ und $\klammer\left\{ \alpha_j, \alpha_k \right\}=\delta_{jk}$ gelten.
Diese Bedingungen können mit $4\times 4$-Matrizen erfüllt werden.
Man verwendet für die Darstellung der Dirac-Gleichung üblicherweise die $\gamma$-Matrizen, welche mit $\beta,\vec{\alpha}$ über
\begin{align}
\gamma^0 &= \beta = \begin{pmatrix}
1 & 0 \\ 0 & -1
\end{pmatrix} &&\tecxt\text{und} & \gamma^j &= \gamma^0\alpha_j = \begin{pmatrix}
0 & \sigma^j \\ -\sigma^j & 0
\end{pmatrix}
\end{align}
in Verbindung gebracht werden können.
Sei $\psi$ nun ein Dirac-Spinorfeld mit $\psi\rightarrow\psi ' = S(\Lambda) \psi$.
Dann ist die Dirac-Gleichung $(\i\slashed{\partial} - m)\psi=0$ Lorentz-invariant, da
\begin{align}
\slashed{\partial}' &= \partial_\mu'\gamma'^\mu = \partial_\mu'S^{-1}\gamma^\mu S = \partial_\mu'\Lambda^\mu_{\ \nu}\gamma^\nu = \slashed{\partial}.
\end{align}
Die Lorentz-Transformation kann hier als
\begin{align}
\psi &\rightarrow\psi' = \kla\left( 1 - \frac{\i}{2}\omega_{\mu\nu}\hat{J}^{\mu\nu} \right) 
\end{align}
dargestellt werden, wobei $\omega_{\mu\nu} = -\omega_{\nu\mu}$ infinitesimale Boosts bzw. Drehungen beschreibt.
Zudem ist
\begin{align}
\hat{J}^{\mu\nu} &= \hat{L}^{\mu\nu} + \hat{S}^{\mu\nu} = \i\kla\left( x^\mu\partial^\nu - x^\nu\partial^\mu  \right) + \frac{\i}{4}\klam\left[ \gamma^\mu, \gamma^\nu \right].
\end{align}
%Insbesondere ist \psi'(x) &= S(\Lambda)\psi(\Lambda^{-1}) = \kla\left( 1 - \frac{\i}{2}\omega_{\mu\nu}S^{\mu\nu} \right) \psi(x^\mu - \omega^{\mu\nu} x_\nu)

\subsection{Lösungen der Dirac-Gleichung}
Es existieren freie Lösungen gegeben durch ebene Wellen
\begin{align}
\psi(x) &= u(p) \e^{\mp\i p\cdot x},
\end{align}
wobei die Koeffizienten die algebraische Gleichung
\begin{align}
0 &= \kla\left( \pm\slashed{p} - m \right)u(p)
\end{align}
erfüllen. 
Die Eigenwerte von $\slashed{p}$ kann man dabei elegant über
\begin{align}
\slashed{p}^2 &= p_\mu p_\nu \gamma^\mu\gamma^\nu = p_\mu p_\nu\kla\left( \klammer\left\{ \gamma^\mu, \gamma^\nu \right\} - \gamma^\nu\gamma^\mu  \right) = p_\mu p_\nu \kla\left( 2g^{\mu\nu} \right) - \slashed{p}^2
\end{align}
bestimmen, da hieraus direkt $\slashed{p}^2 = p^2 = E^2 - \vec{p}\,^2 = m^2$ folgt.
Die Eigenwerte sind also $\pm m$.
Für masselose Teilchen sind diese also trivial.
Wir definieren
\begin{align}
\kla\left( \slashed{p} - m \right)u(p,s) &:= 0 &&\tecxt\text{und}  & \kla\left( \slashed{p} + m \right)v(p,s) &:= 0,
\end{align}
wobei das Argument $s=\pm\frac{1}{2}$ den Spin der Wellenfunkton beschreibt.
Um diese Funktionen interpretieren zu können, wählen wir zunächst $p^\mu = (E, 0, 0 ,p_z)^\top$.
Dann ist
\begin{align}
\slashed{p} &= E\gamma^0 - p_z\gamma^3 = \begin{pmatrix}
E & -p_z\sigma^3 \\ p_z\sigma^3 & -E
\end{pmatrix}.
\end{align}
Insbesondere gilt hier $\klam\left[ \slashed{p}, S_z \right] = 0$, da $S_z = \frac{1}{2}\begin{pmatrix}
\sigma^3 & 0 \\ 0 & \sigma^3
\end{pmatrix}$.
Also haben $\slashed{p}$ und $S_z$ simultane Eigenzustände.
Allgemein findet man $\klam\left[ \slashed{p}, \hat{h} \right] = 0$, wobei $\hat{h} := \frac{\vec{p}\cdot\vec{S}}{|\vec{p}|}$ der Helizitätsoperator ist.
Wir finden in unserem Spezialfall letztlich die vier Lösungen
\begin{align}
u\kla\left( p, \frac{1}{2} \right) &= N\begin{pmatrix}
E + m \\ 0 \\ p_z \\ 0
\end{pmatrix}, &&& u\kla\left( p, -\frac{1}{2} \right) &= N\begin{pmatrix}
0 \\ E + m \\ 0 \\ -p_z
\end{pmatrix}, \\
v\kla\left( p, \frac{1}{2} \right) &= N\begin{pmatrix}
p_z \\ 0 \\ E + m \\ 0
\end{pmatrix} &&\tecxt\text{und}& v\kla\left( p, -\frac{1}{2} \right) &= N\begin{pmatrix}
0 \\ -p_z \\ 0 \\ E + m
\end{pmatrix}.
\end{align}
Die Normierung ist $N=\frac{1}{\sqrt{E+m}}$ und führt mit $u_s(p) := u(p,s)$ auf $\overline{u}_{s_1}u_{s_2} = 2m\delta_{s_1s_2}$.
Analog findet man $\overline{v}_{s_1}v_{s_2} = -2m\delta_{s_1s_2}$. 
\\
Für beliebige $\vec{p}$ kann man diese Lösungen mittels
\begin{align}
u_s'(p) &= \e^{-\frac{\i}{2}\omega_{\mu\nu}S^{\mu\nu}}u_s(p) 
\end{align}
transformieren.
Die $v_s(p)$ Lösungen haben negative Energie, können also wieder als Antiteilchen interpretiert werden.
Allerdings ist hier im Gegensatz zu Klein-Gordon trotzdem $\rho\ge 0$.

\subsubsection{Kopplung ans elektromagnetische Feld}
Wir kennen bereits die Eichtransformation $A'^\mu = A^\mu + \partial^\mu\chi$.
Die Lagrange-Funktion
\begin{align}
L &= \frac{m}{2}\frac{\tecxt\text{d}x^\mu }{\tecxt\text{d}\tau} \frac{\tecxt\text{d}x_\mu }{\tecxt\text{d}\tau} + Q\frac{\tecxt\text{d}x_\mu }{\tecxt\text{d}\tau} A^\mu \kla\left( x(\tau) \right)
\end{align}
führt auf den kanonischen Impuls
\begin{align}
\pi^\mu := \frac{\partial L}{\partial \frac{\tecxt\text{d}x_\mu }{\tecxt\text{d}\tau}} = m\frac{\tecxt\text{d}x^\mu }{\tecxt\text{d}\tau} + Q A^\mu = p^\mu + QA^\mu.
\end{align}
Wir transformieren unseren Impuls also über $p^\mu \rightarrow p^\mu - QA^\mu$ um den zusätzlichen Term zu kompensieren.
Entsprechend führen wir dann mit $p^\mu \equiv \i\partial^\mu$ den Differentialoperator
\begin{align}
\i D^\mu &:= \i\partial^\mu - QA^\mu 
\end{align}
ein.
Dann ist die modifizierte Dirac-Gleichung
\begin{align}
\kla\left( \i \slashed{D} - Q\slashed{A} - m \right)\psi &= 0
\end{align}
eichinvariant, da 
\begin{align}
\slashed{D}'\psi' &= \kla\left( \slashed{D} + \i Q\slashed{\partial}\chi \right) \e^{-\i Q\chi}\psi = \e^{-\i Q\chi} \kla\left( \slashed{D} + \i Q\slashed{\partial}\chi - \i Q\slashed{\partial}\chi \right)\psi = \e^{-\i Q\chi}\slashed{D}\psi.
\end{align}
Man findet die Stromdichte $j^\mu := -Q\overline{\psi}\gamma^\mu \psi$.
Für diese gilt
\begin{align}
D_\mu j^\mu &= -Q\overline{\psi}\slashed{D}\psi - Q\kla\left( \slashed{D}\overline{\psi} \right)\psi = \i Q\overline{\psi}\kla\left( m \right)\psi + \i Q\overline{\psi}\kla\left( -m  \right)\psi =  0.
\end{align}
Die Schrödinger-Gleichung für die minimale Kopplung lautet
\begin{align}
\kla\left( \i\partial_t - Q\phi \right)\psi &= \frac{1}{2m} \kla\left( \vec{p} - Q\vec{A} \right)^2\psi.
\end{align}
Um den nichtrelativistischen Limes der Dirac-Gleichung zu betrachten führen wir zunächst Spinoren $\psi = (\psi_A, \psi_B)^\top$ ein.
Mit $\i D_0 \equiv E - Q\phi$ und $\i D_j \equiv \vec{p} - Q\vec{A} = \vec{\pi}$ können wir dann zunächst
\begin{align}
0 &= \kla\left( \i\slashed{D} - m \right)\psi = \begin{pmatrix}
E - Q\phi - m & -\vec{\sigma}\cdot\kla\vec{\pi} \\ 
\vec{\sigma}\cdot\vec{\pi} & -E + Q\phi - m
\end{pmatrix}\begin{pmatrix}
\psi_A \\ \psi_B
\end{pmatrix}
\end{align}
schreiben. 
Wir nehmen an, dass $\varepsilon := E-m-Q\phi \ll m$.
Dann gilt $\varepsilon \psi_A = \vec{\sigma}\cdot\vec{\pi}\psi_B$ und somit
\begin{align}
0 &= \vec{\sigma}\cdot\vec{\pi}\psi_A - \kla\left( 2m + \varepsilon \right)\psi_B = \frac{1}{\varepsilon}\kla\left( \vec{\sigma}\cdot\vec{\pi} \right)^2\psi_B - \kla\left( 2m + \varepsilon \right)\psi_B.
\end{align}
Mit $\sigma_j\sigma_k = \delta_{jk} + \i\epsilon_{jkl}\sigma_l$ erhalten wir
\begin{align}
\kla\left( \vec{\sigma}\cdot\vec{\pi} \right)^2 &= \sigma_j\sigma_k \kla\left( p_j - QA_j \right)\kla\left( p_k - QA_k \right) = \\
&= \vec{\pi}^2 + \i\epsilon_{jkl}\sigma_l \kla\left( p_jp_k + Q^2A_jA_k  - Qp_jA_k - QA_jp_k \right) = \\
&= \vec{\pi}^2 - \i Q\vec{\sigma}\cdot\kla\left( \vec{p}\times\vec{A} \right) = \vec{\pi}^2 - Q\vec{\sigma}\cdot\vec{B}.
\end{align} 
Daraus folgt dann unter Vernachlässigung von Termen $\mathcal{O}(\varepsilon^2)$ die Pauli-Gleichung
\begin{align}
\varepsilon\psi_B &= \frac{1}{2m}\kla\left( \vec{\pi}^2 - Q\vec{\sigma}\cdot\vec{B} \right)\psi_B.
\end{align}
Verglichen mit der Schrödinger-Gleichung von oben enthält diese einen zusätzlichen Term $\vec{S}\cdot\vec{B}$-Term, der an das Magnetfeld koppelt.
Wir können also auch $\hat{H} = \hat{H}_\tecxt\text{SG} - \vec{\mu}\cdot \vec{B}$ schreiben, wobei das magnetische Moment durch
\begin{align}
\vec{\mu} &:= g_s \frac{Q}{2m}\vec{S}
\end{align}
gegeben ist.
Da $\vec{S} = \frac{1}{2}\begin{pmatrix}
\vec{\sigma}  & 0 \\ 0 & \vec{\sigma}
\end{pmatrix}$, finden wir im Vergleich zur Pauli-Gleichung den gyromagnetischen Faktor $g_s = 2$.
Klassisch war $\vec{\mu} = \frac{Q}{2m}\vec{L}$, also $g=1$. 
Der exakte Wert weicht leicht von 2 ab, was allerdings erst mit QED verstanden werden kann.
Für ein konstantes Magnetfeld können wir $\vec{A} = \frac{1}{2}\vec{B}\times\vec{r}$ verwenden, da 
\begin{align}
\nabla\times\vec{A} &= \frac{1}{2}\vec{B}\kla\left( \nabla\cdot\vec{r} \right) - \frac{1}{2}\kla\left( \vec{B}\cdot\nabla \right)\vec{r} = \frac{3}{2}\vec{B} - \frac{1}{2}\vec{B} = \vec{B}.
\end{align}
Insbesondere gilt dann
\begin{align}
\klammer\left\{ \vec{p}, \vec{A} \right\} &= \frac{1}{2}\epsilon_{jkl}  B_k \klammer\left\{ p_j, x_l \right\} = B_k\epsilon_{ljk}x_lp_j = \vec{B}\cdot\kla\left( \vec{r}\times\vec{p} \right) = \vec{B}\cdot\vec{L}.
\end{align}
Wir finden also 
\begin{align}
\vec{\pi}^2 &= \vec{p}\,^2 + Q^2\vec{A}^2 - Q\klammer\left\{ \vec{p}, \vec{A} \right\} = \vec{p}\,^2 + \frac{Q^2}{4}\vec{B}\cdot\kla\left( \vec{r}\times \kla\left( \vec{B}\times\vec{r} \right) \right) - Q\vec{B}\cdot\vec{L} = \\
&= \vec{p}\,^2 + \frac{Q^2}{4}\kla\left( r^2B^2 - \vec{r}\kla\left( \vec{B}\cdot\vec{r} \right) \right) - Q\vec{B}\cdot\vec{L}.
\end{align}
Hier stimmt $g_l=1$ also mit dem klassischen Wert überein.
\subsubsection{Foldy-Wouthuysen Transformation}
TODO: höhere Ordnungen von $\epsilon$ in der Herleitung oben. Längeres Kapitel mit ausführlichen Herleitungen.

\section{Ununterscheidbare Teilchen}
\subsection{Unterscheidbare Teilchen}
Die Basis für Teilchen 1 sei $\Ket | n^{(1)} \rangle$, und die für Teilchen 2 sei $\Ket | m^{(2)} \rangle$.
Dann können wir beide Teilchen in einem gemeinsamen Produktraum
\begin{align}
\mathcal{H}_2 &:= \mathcal{H}_1^{(1)} \otimes \mathcal{H}_2^{(1)} = \klammer\left\{ \sum_{nm} c_{nm} \Ket | n^{(1)} \rangle \Ket | m^{(2)} \rangle \right\}
\end{align}
beschreiben.
Dieser enthält insbesondere verschränkte Zustände wie zum Beispiel
\begin{align}
\frac{\Ket | 1^{(1)} \rangle \Ket | 2^{(2)} \rangle - \Ket | 2^{(1)} \rangle \Ket | 1^{(2)} \rangle}{\sqrt{2}},
\end{align}
die nicht als einfache Produktzustände dargestellt werden können.
Zwei ein-Teilchen-Observablen $A_1^{(1)}$ und $B_1^{(2)}$ kommutieren.
Da sie sich zunächst auf unterschiedliche Räume beziehen, ist diese Aussage allerdings mathematisch unsauber.
Wir müssen sie zunächst auf den $\mathcal{H}_2$ fortsetzen, also
\begin{align}
A_2 := A_1^{(1)}\otimes 1^{(2)} &&\tecxt\text{und} & B_2 &:= 1^{(1)} \otimes B_1^{(2)}.
\end{align}
Dann gilt $\klam\left[ A_2, B_2 \right] = 0$.
Das bedeutet letztlich einfach, dass wir beispielsweise den Ort von Teilchen 1 und den  Impuls von Teilchen 2 gleichzeitig scharf messen können.
Manchmal schreibt man auch $A_2^{(1)}$ um explizit anzudeuten, dass ein Operator effektiv nur auf einen bestimmten Teilraum wirkt.
\subsection{Identische/Ununterscheidbare Teilchen}
Um Verwirrung zu vermeiden, werden wir überwiegend den Begriff ``identisch'' verwenden, da er beim Lesen leichter von ``unterscheidbar'' zu trennen ist.
\\
Der bisherige Formalismus reicht für diese Idee nicht aus.
Wir können beispielsweise nicht den Zustand
\begin{align}
\Ket | \dots, \psi, \dots,\phi, \dots \rangle &\equiv \text{``Jeweils ein Teilchen im Zustand }\psi\tecxt\text{ und }\phi\tecxt\text{''}
\end{align}
beschreiben.
Wir können nun den Permutationsoperator $P$ einführen, welcher zwei Einträge in einem Zustand vertauscht.
\begin{align}
P_{jk} \Ket | \dots, \underbrace{\psi}_{j\tecxt\text{-ter Eintrag}}, \dots,\underbrace{\phi}_{k\tecxt\text{-ter Eintrag}}, \dots \rangle &= \pm\Ket | \dots, \phi, \dots,\psi, \dots \rangle.
\end{align}
Hierbei gilt das positive Vorzeichen für Bosonen, und das negative für Fermionen.



%Analog zum Vorgehen bei der Klein-Gordon-Gleichung kann man eine Stromdichte konstruieren.
%Betrachte dazu
%\begin{align}
%0 &= \overline{\psi}\kla\left( \i \slashed{D} - m \right)\psi + \psi\kla\left( \i \slashed{D} + m \right)\overline{\psi} = 
%\end{align}

%\begin{thebibliography}{9}
%\bibitem{example} description, \url{https://www.exmapleurl}, day.\,month~2021.
%\end{thebibliography}

\end{document}